\documentclass[11pt]{article}

\usepackage[a4paper,margin=1in]{geometry}
\usepackage{amsmath}
\usepackage{amssymb}
\usepackage{hyperref}
\usepackage{listings}

\newcommand{\dd}{\mathrm{d}}
\newcommand{\pp}{\mathrm{p}}
\newcommand{\qqbar}{q\bar q}
\newcommand{\mll}{m_{\ell\ell}}
\newcommand{\shat}{\hat s}
\newcommand{\ecm}{E_{\mathrm{cm}}}
\newcommand{\mZ}{m_Z}
\newcommand{\sw}{s_w}
\newcommand{\cw}{c_w}
\newcommand{\pb}{\mathrm{pb}}

\title{The \texttt{drell\_yan\_lo} Numerical Experiment}
\author{Computational HEP}
\date{}

\begin{document}
\maketitle

\section{Purpose of the experiment}

This numerical experiment computes a deliberately simple version of
Drell--Yan production at a proton--proton collider:
\[
  \pp\pp \to Z \to \ell^+\ell^- .
\]
The implementation is intentionally pedagogical.  It does not try to be a
precision Drell--Yan prediction.  Instead, it is designed to expose the three
main ingredients that appear in almost every hadron-collider calculation:
\begin{enumerate}
  \item the proton is not an elementary particle, so the incoming hard partons
        carry only fractions \(x_1\) and \(x_2\) of the proton momenta;
  \item the probability densities for finding those partons are the PDFs;
  \item the short-distance partonic reaction is simple and calculable.
\end{enumerate}
In this experiment the hard reaction is
\[
  q \bar q \to Z ,
\]
and the decay \(Z \to \ell^+\ell^-\) is included only through a branching-ratio
factor.  Therefore the final state is treated as an on-shell object with mass
\(\mll\), not as two resolved leptons with an angular distribution.  This is why
the phase space is a \(2\to1\) phase space.

\section{Hadron-collider kinematics}

Let the two proton beams have centre-of-mass energy \(\ecm\).  The proton
momenta are taken to be massless:
\[
  P_1^\mu = \frac{\ecm}{2}(1,0,0,+1), \qquad
  P_2^\mu = \frac{\ecm}{2}(1,0,0,-1).
\]
The hard collision is not between the two full protons.  It is between two
partons inside the protons.  If the parton from proton 1 carries momentum
fraction \(x_1\), and the parton from proton 2 carries momentum fraction
\(x_2\), then
\[
  p_1^\mu = x_1 P_1^\mu, \qquad p_2^\mu = x_2 P_2^\mu .
\]
The partonic centre-of-mass energy is
\[
  \shat = (p_1+p_2)^2 = x_1 x_2 \ecm^2 .
\]

For a \(2\to1\) reaction producing a final-state object of mass \(\mll\), the
partonic energy is fixed:
\[
  \shat = \mll^2 .
\]
It is useful to define
\[
  \tau = x_1 x_2 = \frac{\mll^2}{\ecm^2}.
\]
The longitudinal boost of the partonic collision is described by the rapidity
of the produced object,
\[
  y = \frac{1}{2}\log\frac{x_1}{x_2}.
\]
Solving these two equations gives
\[
  x_1 = \sqrt{\tau}\,e^{+y}, \qquad
  x_2 = \sqrt{\tau}\,e^{-y}.
\]
Because \(0<x_1<1\) and \(0<x_2<1\), the rapidity is bounded:
\[
  -\log\frac{\ecm}{\mll}
  \le y \le
  +\log\frac{\ecm}{\mll}.
\]
The code does not rederive this mapping itself.  It uses the existing
\texttt{FlatPhaseSpace} generator in \(2\to1\) proton-proton mode:
\[
  \texttt{FlatPhaseSpace([0,0], [m\_ll], beam\_types=(1,1))}.
\]
For a \(2\to1\) final state this generator fixes \(\tau=\mll^2/\ecm^2\), samples
only the rapidity variable, and returns \(x_1\), \(x_2\), and the Jacobian.
This is important: the Drell--Yan experiment is not manually inventing a
Bjorken-\(x\) mapping.  It is using the same phase-space machinery as the other
hadron-collider experiments.

\section{PDFs and partonic luminosities}

A PDF \(f_i(x,\mu_F)\) answers the following question:
\begin{quote}
  At factorisation scale \(\mu_F\), how likely is it to find a parton of type
  \(i\) carrying a fraction \(x\) of the proton momentum?
\end{quote}
The index \(i\) can be a gluon, a quark, or an antiquark.  In Drell--Yan at
leading order, only quark--antiquark initial states contribute:
\[
  q\bar q \to Z .
\]
For a proton-proton collider, either beam can provide the quark and the other
beam can provide the antiquark.  Therefore, for each quark flavour \(q\), the
partonic luminosity is the sum of two beam assignments:
\[
  \mathcal{L}_{q\bar q}(x_1,x_2,\mu_F)
  =
  f_q(x_1,\mu_F) f_{\bar q}(x_2,\mu_F)
  +
  f_{\bar q}(x_1,\mu_F) f_q(x_2,\mu_F).
\]
The code keeps these two terms separately in the output table, because it is
instructive to see that they are generally not equal away from \(y=0\).
At positive rapidity, \(x_1>x_2\); at negative rapidity, \(x_2>x_1\).  Since
quark and antiquark PDFs are different functions, the two beam assignments
probe different regions of the proton structure.

LHAPDF returns \(x f_i(x,\mu_F)\), not \(f_i(x,\mu_F)\).  The experiment
therefore converts it back by dividing by \(x\):
\[
  f_i(x,\mu_F) = \frac{x f_i(x,\mu_F)}{x}.
\]
This division is explicit in the code because it is one of the easiest places
to make a factor-of-\(x\) mistake.

\section{The \(Z\)-boson couplings}

The neutral-current \(Z f\bar f\) vertex is written in the convention
\[
  i\frac{e}{2\sw\cw}\gamma^\mu(v_f-a_f\gamma_5).
\]
Here \(e\) is the electromagnetic coupling, and
\[
  \sw = \sin\theta_w, \qquad \cw=\cos\theta_w
\]
are the sine and cosine of the weak mixing angle.  The dimensionless vector
and axial couplings are
\[
  v_f = T^3_f - 2 Q_f \sw^2,
  \qquad
  a_f = T^3_f .
\]
\(Q_f\) is the electric charge in units of the proton charge.  \(T^3_f\) is the
third component of weak isospin.  For example,
\[
  (Q_u,T^3_u)=\left(+\frac{2}{3},+\frac{1}{2}\right),
  \qquad
  (Q_d,T^3_d)=\left(-\frac{1}{3},-\frac{1}{2}\right).
\]
The code stores these charge and isospin assignments for \(d,u,s,c,b\) quarks.
For a chosen flavour \(q\), it computes \(v_q\) and \(a_q\), and then weights the
channel by \(v_q^2+a_q^2\).

\section{The hard matrix element}

The partonic process is
\[
  q(p_1) + \bar q(p_2) \to Z(k),
  \qquad k^2=\mll^2 .
\]
The amplitude is
\[
  \mathcal{M}
  =
  \bar v(p_2)
  \left[
    i\frac{e}{2\sw\cw}\gamma^\mu(v_q-a_q\gamma_5)
    \right]
  u(p_1)\,
  \epsilon_\mu^*(k).
\]
The experiment uses the spin- and colour-averaged squared matrix element
\[
  \overline{|\mathcal{M}|^2}
  =
  \frac{g_Z^2\,\shat}{12}\left(v_q^2+a_q^2\right),
  \qquad
  g_Z = \frac{e}{\sw\cw}.
\]
Let us unpack this formula.
\begin{itemize}
  \item The bar over \(|\mathcal{M}|^2\) means that we average over the
        unobserved spin and colour states of the incoming quark and antiquark.
  \item The factor \(1/4\) from spin averaging is already included.
  \item The factor \(1/3\) from colour averaging is already included.  The
        final \(Z\) is colourless, so only a colour-singlet quark-antiquark
        combination contributes.
  \item The interference between vector and axial terms vanishes after the
        spin sum for this fully inclusive \(2\to1\) production rate.
\end{itemize}
This is what is meant by ``the LO matrix element is trivial'' in this numerical
experiment: the code does not call a generated matrix-element library.  It uses
this one analytic expression directly.

\section{The \(2\to1\) partonic cross section}

For a \(2\to1\) process, the partonic cross section contains a delta function
that enforces the on-shell condition:
\[
  \dd \hat\sigma_q
  =
  \frac{\pi}{\shat}
  \overline{|\mathcal{M}_q|^2}
  \delta(\shat-\mll^2).
\]
This expression is compact, so it is worth saying what it means.
The incoming partons can have many possible \(x_1,x_2\), but only combinations
with
\[
  x_1 x_2 \ecm^2 = \mll^2
\]
can produce an on-shell object of mass \(\mll\).  The delta function removes
one integration variable.  In the code, the phase-space generator handles this
for us.  In its \(2\to1\) mode it fixes \(\tau=\mll^2/\ecm^2\) and leaves only
the rapidity integral.

The hadronic cross section in this approximation is therefore
\[
  \sigma(\pp\pp\to Z\to\ell^+\ell^-)
  =
  \sum_q
  \int \dd x_1\,\dd x_2\,
  \mathcal{L}_{q\bar q}(x_1,x_2,\mu_F)\,
  \frac{\pi}{\shat}
  \overline{|\mathcal{M}_q|^2}
  \delta(\shat-\mll^2)\,
  \mathrm{BR}(Z\to\ell^+\ell^-).
\]
Equivalently, after using \(\tau\) and \(y\),
\[
  \sigma
  =
  \sum_q
  \int_{-Y}^{+Y}\dd y\,
  \frac{1}{\ecm^2}\,
  \mathcal{L}_{q\bar q}(x_1(y),x_2(y),\mu_F)\,
  \pi\,
  \frac{\overline{|\mathcal{M}_q|^2}}{\mll^2}\,
  \mathrm{BR}(Z\to\ell^+\ell^-),
\]
where
\[
  Y = \log\frac{\ecm}{\mll}.
\]
The factor \(1/\ecm^2\) is the Jacobian from the delta function
\(\delta(x_1x_2\ecm^2-\mll^2)\).  In the code this is part of the
\texttt{FlatPhaseSpace} Jacobian.

\section{The branching ratio}

By default the experiment returns a cross section for
\[
  \pp\pp\to Z\to \ell^+\ell^- .
\]
The branching ratio is estimated at tree level.  For one charged lepton flavour,
\[
  \Gamma(Z\to\ell^+\ell^-)
  =
  \frac{G_F\mZ^3}{6\sqrt{2}\pi}
  \left(v_\ell^2+a_\ell^2\right),
\]
and
\[
  \mathrm{BR}(Z\to\ell^+\ell^-)
  =
  \frac{\Gamma(Z\to\ell^+\ell^-)}{\Gamma_Z}.
\]
Here \(\Gamma_Z\) is taken from the model parameters.  If one wants inclusive
on-shell \(Z\) production rather than a charged-lepton final state, one can run
with
\[
  \texttt{--branching\_ratio 1}.
\]

\section{Monte Carlo integration}

The integrator samples a point \(u\in[0,1]\).  The phase-space generator maps
this to a rapidity
\[
  y(u)= -Y + 2Yu,
\]
and then to
\[
  x_1(u)=\sqrt{\tau}\,e^{y(u)}, \qquad
  x_2(u)=\sqrt{\tau}\,e^{-y(u)}.
\]
At that point the code computes
\[
  w(u)
  =
  J(u)
  \sum_q
  \left[
    f_q(x_1,\mu_F)f_{\bar q}(x_2,\mu_F)
    +
    f_{\bar q}(x_1,\mu_F)f_q(x_2,\mu_F)
    \right]
  \pi
  \frac{\overline{|\mathcal{M}_q|^2}}{\mll^2}
  \mathrm{BR}(Z\to\ell^+\ell^-),
\]
where \(J(u)\) is the phase-space/Bjorken-\(x\) Jacobian returned by the
generator.  The unit conversion
\[
  1~\mathrm{GeV}^{-2} = 0.389379338\times 10^9~\pb
\]
is applied at the end.

The integral over \(u\) gives the total cross section:
\[
  \sigma = \int_0^1 \dd u\, w(u).
\]
The numerical integrator reports an estimate and an integration error after
each iteration.

\section{Rapidity distribution}

In addition to the Monte Carlo integral, the experiment writes a simple
rapidity table.  This table is evaluated deterministically at bin centres.
Since the generator weight is a weight per unit \(u\), and
\[
  \dd y = 2Y\,\dd u,
\]
the code divides by the rapidity range \(2Y\) to obtain
\[
  \frac{\dd\sigma}{\dd y}.
\]
The output files are
\[
  \texttt{drell\_yan\_lo\_rapidity.csv},
  \qquad
  \texttt{drell\_yan\_lo\_rapidity.png}.
\]
The CSV table also includes separate channel contributions such as
\(\texttt{u qbar}\) and \(\texttt{ubar q}\).  These labels indicate which beam
provided the quark and which beam provided the antiquark.

\section{Example commands}

A default run near the \(Z\) pole is
\begin{lstlisting}[language=bash]
./run.py drell_yan_lo -ni 10 -npi 1000 \
  --e_cm 13000 \
  --quark_flavours d,u,s,c \
  --lepton mu \
  --output_dir drell_yan_lo_output
\end{lstlisting}
For inclusive on-shell \(Z\) production, omit the branching-ratio suppression:
\begin{lstlisting}[language=bash]
./run.py drell_yan_lo -ni 10 -npi 1000 \
  --branching_ratio 1 \
  --output_dir drell_yan_lo_inclusive
\end{lstlisting}

\section{What is deliberately not included}

This experiment is intentionally minimal.  The following physics effects are
not included:
\begin{itemize}
  \item photon exchange, \(\gamma^*\to\ell^+\ell^-\);
  \item \(\gamma\)-\(Z\) interference;
  \item the finite-width Breit-Wigner line shape of the \(Z\);
  \item the explicit two-lepton decay angles;
  \item QCD radiation such as \(q\bar q\to Zg\) or \(qg\to Zq\);
  \item virtual loop corrections;
  \item detector cuts on the charged leptons.
\end{itemize}
These omissions are not accidents.  The purpose of the experiment is to show,
as cleanly as possible, how a hadron-collider observable is built from
partonic luminosities and a hard matrix element.

\end{document}
